\documentclass[UTF8,11pt]{ctexart}
\usepackage[a4paper,margin=24mm]{geometry}
\usepackage{amsmath,amssymb,amsthm,mathtools,mathrsfs}
\usepackage[all]{xy}
\usepackage{xurl,hyperref,enumitem}
\usepackage{tikz}
\usetikzlibrary{arrows.meta,cd,decorations.markings,calc,patterns}
\hypersetup{hidelinks,breaklinks=true}
\setlength{\emergencystretch}{3em}
\allowdisplaybreaks
\newtheorem*{theorem}{Theorem}
\newtheorem*{thm}{Theorem}
\newtheorem*{lemma}{Lemma}
\newtheorem*{proposition}{Proposition}
\newtheorem*{prop}{Proposition}
\newtheorem*{corollary}{Corollary}
\newtheorem*{cor}{Corollary}
\newtheorem*{claim}{Claim}
\theoremstyle{definition}
\newtheorem*{definition}{Definition}
\newtheorem*{defn}{Definition}
\newtheorem*{remark}{Remark}
\newtheorem*{example}{Example}
\newcommand{\R}{\mathbb R}
\newcommand{\C}{\mathbb C}
\newcommand{\Z}{\mathbb Z}
\newcommand{\Q}{\mathbb Q}
\newcommand{\N}{\mathbb N}
\newcommand{\F}{\mathbb F}
\newcommand{\D}{\mathbb D}
\newcommand{\bB}{\mathbb B}
\newcommand{\bC}{\mathbb C}
\newcommand{\bR}{\mathbb R}
\newcommand{\bZ}{\mathbb Z}
\newcommand{\bN}{\mathbb N}
\newcommand{\bQ}{\mathbb Q}
\newcommand{\bD}{\mathbb D}
\newcommand{\bF}{\mathbb F}
\newcommand{\CP}{\mathbb{CP}}
\newcommand{\RP}{\mathbb{RP}}
\newcommand{\sO}{\mathscr O}
\newcommand{\sI}{\mathscr I}
\newcommand{\sF}{\mathscr F}
\newcommand{\sG}{\mathscr G}
\newcommand{\sH}{\mathscr H}
\newcommand{\sR}{\mathscr R}
\newcommand{\sM}{\mathscr M}
\newcommand{\sV}{\mathscr V}
\newcommand{\sU}{\mathscr U}
\DeclareMathOperator{\Hom}{Hom}
\DeclareMathOperator{\Ext}{Ext}
\DeclareMathOperator{\Tor}{Tor}
\DeclareMathOperator{\Spec}{Spec}
\DeclareMathOperator{\Proj}{Proj}
\DeclareMathOperator{\supp}{supp}
\DeclareMathOperator{\im}{im}
\DeclareMathOperator{\id}{id}
\DeclareMathOperator{\Tr}{Tr}
\DeclareMathOperator{\tr}{tr}
\DeclareMathOperator{\rank}{rank}
\DeclareMathOperator{\ord}{ord}
\DeclareMathOperator{\dist}{dist}
\DeclareMathOperator{\End}{End}
\DeclareMathOperator{\Aut}{Aut}
\DeclareMathOperator{\Gal}{Gal}
\DeclareMathOperator{\vol}{vol}
\DeclareMathOperator{\Cl}{Cl}
\DeclareMathOperator{\coker}{coker}
\DeclareMathOperator{\codim}{codim}
\DeclareMathOperator{\divergence}{div}
\DeclareMathOperator{\Ric}{Ric}
\DeclareMathOperator{\grad}{grad}
\DeclareMathOperator{\Hess}{Hess}
\newcommand{\abs}[1]{\left\lvert#1\right\rvert}
\newcommand{\sabs}[1]{\lvert#1\rvert}
\newcommand{\babs}[1]{\bigl\lvert#1\bigr\rvert}
\newcommand{\Babs}[1]{\Bigl\lvert#1\Bigr\rvert}
\newcommand{\bbabs}[1]{\biggl\lvert#1\biggr\rvert}
\newcommand{\BBabs}[1]{\Biggl\lvert#1\Biggr\rvert}
\newcommand{\norm}[1]{\left\lVert#1\right\rVert}
\newcommand{\snorm}[1]{\lVert#1\rVert}
\newcommand{\bnorm}[1]{\bigl\lVert#1\bigr\rVert}
\newcommand{\Bnorm}[1]{\Bigl\lVert#1\Bigr\rVert}
\newcommand{\bbnorm}[1]{\biggl\lVert#1\biggr\rVert}
\newcommand{\BBnorm}[1]{\Biggl\lVert#1\Biggr\rVert}
\newcommand{\widebar}[1]{\overline{#1}}
\newcommand{\nicefrac}[2]{\frac{#1}{#2}}
\newcommand{\linnprod}[2]{\langle#1,#2\rangle}
\newcommand{\myindex}[1]{#1}
\newcommand{\glsadd}[1]{}
\newcommand{\avoidbreak}{}
\newcommand{\sourceref}[1]{\textnormal{[原文：\nolinkurl{#1}]}}
\newcommand{\thmref}[1]{\sourceref{#1}}
\newcommand{\lemmaref}[1]{\sourceref{#1}}
\newcommand{\propref}[1]{\sourceref{#1}}
\newcommand{\corref}[1]{\sourceref{#1}}
\newcommand{\defnref}[1]{\sourceref{#1}}
\newcommand{\exampleref}[1]{\sourceref{#1}}
\newcommand{\exerciseref}[1]{\sourceref{#1}}
\newcommand{\figureref}[1]{\sourceref{#1}}
\newcommand{\sectionref}[1]{\sourceref{#1}}
\newcommand{\appendixref}[1]{\sourceref{#1}}
\newcommand{\stacksref}[2]{\href{https://stacks.math.columbia.edu/tag/#1}{#2 [#1]}}

\begin{document}
\section*{034\quad 全纯函数芽环的Noetherian有限生成性}
\subsection*{来源与计数目标}
Ji\v{r}\'i Lebl, \emph{Tasty Bits of Several Complex Variables}。从作者公开的原生 TeX 正文摘录；获取日期：2026-09-28。使用实际访问的在线正文，不把工作分支冒称冻结的印刷版。
\par\url{https://raw.githubusercontent.com/jirilebl/scv/master/scv.tex}
\par\url{https://www.jirka.org/scv/}
原文位置：Chapter 6，Weierstrass division theorem 及随后环结构节的 $\sO_p$ is Noetherian 定理。
\par\textbf{本文件只计一个问题：}全纯函数芽环中每个理想有限生成；除法构造仅为支撑，不另计题数。
\subsection*{作者命题与人类证明}

\paragraph{Weierstrass division: the existence construction used below.}
\begin{thm}[Weierstrass division theorem, original statement]
Suppose $f$ is holomorphic near the origin, and suppose $P$ is a Weierstrass polynomial of degree $k\geq1$ in $z_n$. Then there exists a neighborhood $V$ of the origin and unique $q,r\in\sO(V)$, where $r$ is a polynomial in $z_n$ of degree less than $k$, and on $V$,
\[f=qP+r.\]
\end{thm}
Note that $r$ need not be a Weierstrass polynomial; it need not be monic nor do the coefficients need to vanish at the origin. It is simply a polynomial in $z_n$ with coefficients that are holomorphic functions of the first $n-1$ variables.
\begin{proof}
Uniqueness is left as an exercise. There exists a connected neighborhood $V=V'\times D$ of the origin, where $D$ is a disc, $f$ and $P$ are continuous on $V'\times\widebar D$, and $P$ is not zero on $V'\times\partial D$. Let
\[
q(z',z_n)=\frac1{2\pi i}\int_{\partial D}\frac{f(z',\zeta)}{P(z',\zeta)(\zeta-z_n)}\,d\zeta.
\]
As $P$ is not zero on $V'\times\partial D$, the function $q$ is holomorphic in $V$ (differentiate under the integral). If $P$ did divide $f$, then $q$ would really be $\frac fP$. But if $P$ does not divide $f$, then the Cauchy integral formula does not apply and $q$ is not equal to $\frac fP$. Interestingly, the expression does give the quotient in the division with remainder. Write $f$ using the Cauchy integral formula in $z_n$ and subtract $qP$ to obtain $r$:
\begin{align*}
r(z',z_n)&=f(z',z_n)-q(z',z_n)P(z',z_n)\\
&=\frac1{2\pi i}\int_{\partial D}\frac{f(z',\zeta)P(z',\zeta)-f(z',\zeta)P(z',z_n)}{P(z',\zeta)(\zeta-z_n)}\,d\zeta.
\end{align*}
We need to show $r$ is a polynomial in $z_n$ of degree less than $k$. In the expression inside the integral, the numerator is of the form $\sum_{\ell=1}^k h_\ell(z',\zeta)(\zeta^\ell-z_n^\ell)$ and is therefore divisible by $(\zeta-z_n)$. The numerator is a polynomial of degree $k$ in $z_n$. After dividing by $(\zeta-z_n)$, the integrand becomes a polynomial in $z_n$ of degree $k-1$. Use linearity of the integral to integrate the coefficients of the polynomial. Each coefficient is a holomorphic function in $V'$ and the proof is finished. Some coefficients may have integrated to zero, so we can only say that $r$ is a polynomial of degree $k-1$ or less.
\end{proof}

\paragraph{Noetherian property (the counted problem).}
A commutative ring $R$ is Noetherian if every ideal in $R$ is finitely generated. That is, for every ideal $I\subset R$ there exist finitely many generators $f_1,\ldots,f_k\in I$: Every $g\in I$ can be written as $g=c_1f_1+\cdots+c_kf_k$, for some $c_1,\ldots,c_k\in R$.
\begin{thm}
$\sO_p$ is Noetherian.
\end{thm}
\begin{proof}
Without loss of generality, $p=0$. The proof is by induction on dimension. By \exerciseref{exercise:onedimOisPID}, ${}_1\sO_0$ is Noetherian. By \exerciseref{exercise:biholringiso}, we are allowed a biholomorphic change of coordinates near the origin. For induction, suppose ${}_{n-1}\sO_0$ is Noetherian and let $I\subset{}_n\sO_0$ be an ideal. If $I=\{0\}$ or $I={}_n\sO_0$, then the assertion is obvious. Therefore, assume that all elements of $I$ vanish at the origin ($I\neq{}_n\sO_0$), and that there exist elements that are not identically zero ($I\neq\{0\}$). Let $g$ be such an element. After perhaps a linear change of coordinates, assume $g$ is a Weierstrass polynomial in $z_n$ by the preparation theorem. The ring ${}_{n-1}\sO_0[z_n]$ is a subring of ${}_n\sO_0$. The set $J={}_{n-1}\sO_0[z_n]\cap I$ is an ideal in the ring ${}_{n-1}\sO_0[z_n]$. By the Hilbert basis theorem (see \thmref{thm:hilbertbasis} in the appendix for a proof), as ${}_{n-1}\sO_0$ is Noetherian, the ring ${}_{n-1}\sO_0[z_n]$ is also Noetherian. Thus $J$ has finitely many generators, that is, $J=(h_1,\ldots,h_k)$ in the ring ${}_{n-1}\sO_0[z_n]$. By the division theorem, every $f\in I$ is of the form $f=qg+r$, where $r\in{}_{n-1}\sO_0[z_n]$ and $q\in{}_n\sO_0$. As $f$ and $g$ are in $I$, so is $r$. As $g$ and $r$ are in ${}_{n-1}\sO_0[z_n]$, they are both in $J$. Write $g=c_1h_1+\cdots+c_kh_k$ and $r=d_1h_1+\cdots+d_kh_k$. Then
\[f=(qc_1+d_1)h_1+\cdots+(qc_k+d_k)h_k.\]
So $h_1,\ldots,h_k$ also generate $I$ in ${}_n\sO_0$.
\end{proof}

\subsection*{整理说明（非作者正文）}
记号补明：${}_n\sO_p$ 是 $\C^n$ 在 $p$ 处的全纯函数芽环；Weierstrass 多项式在 $z_n$ 中首一，其他系数为 $z'$ 的全纯芽并在原点消失。标准先修为 Hilbert 基定理、一变量全纯芽按消失阶分解、坐标变换诱导函数芽环同构，以及 Weierstrass 预备定理（本批031已有作者证明）。

原作者将除法定理的唯一性留作练习，本文件如实保留该句；本题 Noetherian 证明只使用除法的存在性，其积分构造与余式次数论证已完整收录。不存在把除法定理当成一份独立完整题证计数的情况。难度依据是把解析理想经除法余式接到低一维的多项式理想，再将有限生成元提升回来，不是仅引用“函数芽环 Noetherian”作为结论。外部引用保留原书标签。
\subsection*{可修改的简短标注（非作者正文）}
$c$：局部解析函数环中理想的有限生成性。

$a$：Weierstrass 预备分解；Cauchy 积分除法；有限次数余式；多项式子环；Hilbert 基定理；维数归纳。

\texttt{cross\_theory\_status = yes}。解析除法将无限展开的函数芽压到多项式子环，使代数的 Hilbert 基定理能够产生原解析理想的有限生成元。位置：除法的积分构造及 Noetherian 证明的 $J$ 与 $r$。

全部标注均为非穷尽、可修改初稿；未标注不表示未使用。
\subsection*{许可与修改记录}
原数学正文作者：Ji\v{r}\'i Lebl。作者网站及源书声明允许 CC BY-SA 4.0 或 CC BY-NC-SA 4.0 双许可；本条选择 CC BY-SA 4.0，整理部分亦采用同一许可。\url{https://creativecommons.org/licenses/by-sa/4.0/}。2026-09-28：助手截取题证、调整独立排版，加入来源、范围与初稿标注；不暗示作者认可本次整理。保留的是所用 TeX 正文摘录，未将其称为整书原件。
\end{document}
